\documentclass[final, finnished, monochromatic]{jyuposter}

\addbibresource{thesis.bib}

\newcommand \jyuauthor{Author Name, Other Author Name}
\newcommand \jyuauthorfinnish{Tekijän Nimi, Toisen Tekijän Nimi}
\newcommand \jyutitle{Title of Work}
\newcommand \jyutitlefinnish{Työn otsikko}
\newcommand \jyupublication{Type of Work}
\newcommand \jyupublicationfinnish{Työn tyyppi}
\newcommand \jyusupervisor{Supervisor Name, Other Supervisor Name}
\newcommand \jyusupervisorfinnish{Ohjaajan Nimi, Toisen Ohjaajan Nimi}
\newcommand \jyuinstitute{University of Jyväskylä}
\newcommand \jyuinstitutefinnish{Jyväskylän yliopisto}
\newcommand \jyufaculty{Faculty of Information Technology}
\newcommand \jyufacultyfinnish{Informaatioteknologian tiedekunta}
\newcommand \jyudepartment{Department of Mathematical Information Technology}
\newcommand \jyudepartmentfinnish{Tietotekniikan laitos}
\newcommand \jyuurl{https://www.jyu.fi/it/en/}
\newcommand \jyuurlfinnish{https://www.jyu.fi/it/fi/}
\newcommand \jyuslogan{Serving the Future since 1863.}
\newcommand \jyusloganfinnish{Tulevaisuuden palveluksesssa jo vuodesta 1863.}
\newcommand \jyushortslogan{Since 1863.}
\newcommand \jyushortsloganfinnish{Vuodesta 1863.}
\newcommand \jyuyear{2017}
\newcommand \jyumonth{05}
\newcommand \jyuday{01}

\makesetup

% These commands are optional.
\newcommand \yesnumber{\addtocounter{equation} 1\tag \theequation}
\newcommand \full{{\mathrm d}}

\begin{document}

\maketitle

\begin{multicols} 3

\noindent
\begin{tcolorbox}
\section*{Abstract}
\fontfamily{fla}\fontseries b\fontshape n\fontsize{18}{27}\selectfont

Yhtälö
\begin{align*}
  \phi (R', R'') & = \int \full x \chi (t, x)^* \psi (t, x)
  \yesnumber \label{e/equation}
\end{align*}
löytyy myöhemmin vielä tekstin seasta numerolla \ref{e/equation-again}.

\lipsum[1]
\end{tcolorbox}

\section*{Piirasta}

\lipsum[1-2]

Logo purkissa.

\renewcommand \jyuhandlestyle{jyublue}
\renewcommand \jyuflamestyle{jyuorange}
% \fbox{\jyulogo{width=50mm}}
% \begin{tikzpicture}[xscale=1, yscale=-1]
\fill[\jyuhandlestyle]
(257.512, 139.681)
.. controls (233.408, 139.681) and (214.211, 125.393) .. (204.656, 110.705)
-- (204.463, 110.409)
-- (310.557, 110.409)
-- (310.373, 110.702)
.. controls (298.939, 128.926) and (279.592, 139.681) .. (257.512, 139.681)
-- cycle;
\fill[\jyuhandlestyle]
(266.711, 240.861)
-- (248.308, 240.861)
-- (248.308, 144.143)
-- (248.549, 144.210)
.. controls (255.668, 146.157) and (259.369, 146.448) .. (266.465, 144.212)
-- (266.711, 144.133)
-- cycle;
\fill[\jyuflamestyle]
(267.338000, 11.798300)
.. controls (276.326000, 20.113200) and (274.309000, 32.512200) .. (272.375000, 36.259700)
.. controls (268.160000, 44.428700) and (258.779000, 54.937000) .. (256.408000, 57.811500)
.. controls (252.453000, 62.606400) and (240.572000, 74.299300) .. (239.287000, 80.762200)
.. controls (238.207000, 86.208000) and (239.916000, 94.199700) .. (240.653000, 97.158600)
.. controls (238.020000, 97.158600) and (233.195000, 97.155700) .. (233.195000, 97.155700)
.. controls (231.297000, 92.242600) and (228.225000, 83.233800) .. (229.989000, 73.049300)
.. controls (231.287000, 65.562400) and (241.873000, 53.722600) .. (243.829000, 51.498000)
.. controls (247.334000, 47.514100) and (257.143000, 35.577600) .. (261.129000, 30.399800)
.. controls (265.672000, 24.501900) and (265.332000, 20.722100) .. (262.277000, 15.980400)
.. controls (262.322000, 16.028300) and (267.338000, 11.798300) .. (267.338000, 11.798300)
-- cycle;
\fill[\jyuflamestyle]
(259.482000, 107.087000)
-- (246.650000, 107.087000)
.. controls (244.024000, 99.354000) and (240.686000, 89.173800) .. (243.857000, 81.142500)
.. controls (247.268000, 72.506300) and (262.770000, 60.119600) .. (272.395000, 49.342200)
.. controls (277.570000, 42.855400) and (279.510000, 38.771400) .. (276.387000, 33.179100)
.. controls (276.387000, 33.179100) and (281.486000, 28.977500) .. (281.486000, 28.961900)
.. controls (286.928000, 33.636200) and (290.895000, 46.058100) .. (284.967000, 53.727500)
.. controls (281.383000, 58.362700) and (277.672000, 62.325100) .. (272.855000, 67.511700)
.. controls (268.160000, 72.570700) and (263.506000, 76.935000) .. (257.158000, 84.233800)
.. controls (252.301000, 89.818800) and (257.730000, 101.007000) .. (259.482000, 107.087000)
-- cycle;
\fill[\jyuflamestyle]
(274.195000, 107.087000)
-- (263.305000, 107.087000)
.. controls (262.279000, 103.203000) and (260.643000, 99.126400) .. (260.008000, 95.158600)
.. controls (259.230000, 90.303700) and (259.564000, 87.538500) .. (263.039000, 83.995100)
.. controls (268.016000, 78.922300) and (280.275000, 70.687900) .. (283.219000, 66.920400)
.. controls (285.758000, 63.670800) and (284.252000, 59.470200) .. (284.252000, 59.470200)
.. controls (284.250000, 59.478000) and (289.289000, 55.278800) .. (289.264000, 55.278800)
.. controls (295.498000, 62.517000) and (293.779000, 73.065400) .. (289.107000, 77.040500)
.. controls (285.713000, 79.925700) and (278.418000, 85.955500) .. (273.613000, 90.771900)
.. controls (271.018000, 93.405700) and (270.846000, 93.600000) .. (271.750000, 97.701600)
.. controls (272.502000, 100.847000) and (273.551000, 104.870000) .. (274.195000, 107.087000)
-- cycle;
\fill[\jyuflamestyle]
(226.463000, 77.370100)
.. controls (223.523000, 73.077100) and (222.094000, 69.014600) .. (221.760000, 63.613700)
.. controls (221.012000, 51.454000) and (229.762000, 41.539000) .. (237.028000, 31.329000)
.. controls (243.256000, 22.577100) and (249.092000, 14.712300) .. (249.795000, 11.304100)
.. controls (250.319000, 8.772900) and (249.273000, 5.995060) .. (248.480000, 4.293400)
.. controls (248.472000, 4.294380) and (253.088000, 0.401314) .. (253.623000, -0.000053)
.. controls (260.328000, 6.125440) and (260.730000, 15.604900) .. (257.334000, 22.763600)
.. controls (249.711000, 38.827100) and (224.246000, 59.649400) .. (226.463000, 77.370100)
-- cycle;
\end{tikzpicture}


\bigskip
\noindent
\includegraphics[width=\columnwidth]{pie}

\lipsum[4-5]

\bigskip
\noindent
\begin{minipage} \columnwidth
\begin{flushright}
\includegraphics[width=\columnwidth]{equipment} \\
\fontfamily{fla}\fontseries m\fontshape{it}\fontsize{18}{22}\selectfont
Tässä on kaikenlaista kuvatekstiä.
\end{flushright}
\end{minipage}

\lipsum[6-7]

\section*{Kvanttipelleilyä}

Tätä asiakirjaluokkaa voi käyttää vaikka
kvanttimekaniikkaan liittyvän \cite{feynman-1948} yhtälön
\begin{align*}
  \phi (R', R'') & = \int \full x \chi (t, x)^* \psi (t, x)
  \yesnumber \label{e/equation-again}
\end{align*}
kirjoittamiseen.
On tosin parempi olla tekemättä niin.

\bigskip
\noindent
\includegraphics[width=\columnwidth]{bars}

\bigskip
\noindent
\begin{tcolorbox}
\section*{Conclusions}

\begin{itemize}
  \item Ergh!
  \item Urgh!
\end{itemize}
\end{tcolorbox}

\nocite{*}

\printbibliography

\end{multicols}

\end{document}
